\documentclass[aps,prb,onecolumn,nofootinbib,citeautoscript,10pt]{revtex4-2} 


\synctex=1 


\usepackage{amsmath,amssymb} 
\usepackage{comment}

\usepackage[tight]{subfigure} 

\usepackage[dvipsnames]{xcolor} 
\usepackage[papersize={8.5in,11in}]{geometry}
\usepackage[colorlinks=true]{hyperref}
\hypersetup{
    bookmarks=true,         % show bookmarks bar?
    unicode=false,          % non-Latin characters 
    pdftoolbar=true,        % show Acrobat
    pdfmenubar=true,        % show Acrobat 
    pdffitwindow=false,     % window fit to page when opened
    pdfstartview={FitH},    % fits the width of the page to the window
    pdfkeywords={keyword1} {key2} {key3}, % list of keywords
    pdfnewwindow=true,      % links in new window
    colorlinks=true,       % false: boxed links; true: colored links
    linkcolor=magenta, %red,          % color of internal links (change box color with linkbordercolor)
    citecolor=blue,        % color of links to bibliography
    filecolor=magenta,      % color of file links
    urlcolor=blue           % color of external links
} 

\geometry{top=1.2 cm, left= 1.2 cm, right= 1.2 cm, bottom= 1.2 cm}        



%========================= new packages for deletion ==================================
\usepackage[normalem]{ulem}
\usepackage{xcolor}

\newcommand{\cdel}[1]{\textcolor{red}{\sout{#1}}}
\newcommand{\cadd}[1]{\textcolor{blue}{#1}}
\newcommand{\eqn}[1]{\textcolor{red}{(#1)}}
\DeclareMathOperator{\arctanh}{arctanh}



%========================= PACKAGES ==================================
\usepackage{dcolumn}
\usepackage{color}
\usepackage{amssymb,amsmath}
\usepackage{tabularx,graphicx}
\usepackage{epstopdf}
\usepackage{latexsym}
\usepackage{colortbl}
\usepackage{psfrag}
\usepackage{bbm,bm,array,physics}
\usepackage{dsfont}
\usepackage{float, mathrsfs}

\def \nn{\nonumber \\}
\def\*#1{\boldsymbol{#1}} %\bold font
\newcommand{\at}[2][]{#1|_{#2}}
\DeclareRobustCommand{\bs}[1]{\bm{#1}}

\newcommand{\figuretag}[1]{%
  \addtocounter{figure}{-1}%
  \renewcommand{\thefigure}{#1}%
}

%=======================================================================
%========================= ARTICLE STARTS HERE =========================
%=======================================================================
\begin{document}

\noindent Dear Editor,\\
\\ \noindent Thank you very much for handling our manuscript and for the constructive reports of the two Reviewers. We have addressed each of the comments below, and the comments are highlighted in blue. All the corresponding changes are marked in red in the revised manuscript. We trust that the manuscript is now suitable for publication.
\\\\ \noindent Sincerely,

\noindent The Authors

\vspace{0.5 cm}

\section{Response to Reviewer \#1}

\begin{enumerate}

%%%%%%%%%%%%%%%%%%%%%%%%%
\item \textcolor{blue}{Validity of the strain-induced Hamiltonian and its derivation. The justification for applying the results derived for untilted, isotropic Dirac cones (Ref.~[5]) to tilted, anisotropic cones is insufficient. The assumption ``we assume that this applies also to the case of the tilted versions'' is non-trivial and requires rigorous justification, since tilt and anisotropy can alter how strain couples to the Dirac fermions, particularly when the Dirac point is shifted away from $K$/$K^\prime$. The authors must provide a first-principles derivation or cite literature that rigorously extends the formalism. Furthermore, the claim that ``the strain [is] applied in such a way that it affects mainly the off-diagonal terms arising from the nn hoppings'' is vague. What specific strain configuration ensures this, and what are the constraints on the magnitude and orientation of the strain?}

We thank the Referee for this comment, which concerns the foundation of the calculation. We agree that the extension of the strain coupling, derived for untilted isotropic cones at $K$ and $K^\prime$, to tilted and anisotropic cones is an assumption and not a result. A first-principles derivation requires a specific lattice model, with its nnn hoppings and the location of the Dirac points, and we have not carried it out. We have therefore made three changes. First, the sentence that presented the extension as justified has been replaced by one stating that it is motivated by the observation that the Dirac point and the $K$ point of the deformed lattice do not coincide in general. Second, the assumptions are now stated explicitly. The tilt, which originates from the nnn hoppings, enters only the term proportional to the identity matrix, while the strain acts on the nn hoppings and enters the off-diagonal part. The coefficients of the terms of order $a_0\,k$ are those obtained by expanding around $K$, and they are modified for a Dirac point displaced from $K$. We state that these coefficients are the ones of the model. Third, we have specified the strain configuration, which is a strain $\varsigma_{yy}$ along the bond direction $\bs M_1$ with a constant gradient along the same direction, and we have derived the conditions on its magnitude. They are $|\tilde\gamma|\ll1$ and $a_0\,|k_x|\ll1$, and for a strained region of width $W$ they translate into a window for the pseudomagnetic field. The new text, placed after the definition of the Hamiltonian in Sec.~II~A, reads as follows.\\``We make the assumptions behind the Hamiltonian of Sec.~II~A explicit. The tilt originates from the nnn hoppings, which connect sites of the same sublattice and therefore contribute, at linear order in the momentum, only to the term proportional to the identity matrix. The strain-induced terms originate from the modulation of the nn hoppings, which connect the two sublattices and therefore enter the off-diagonal part $\mathcal H_s$. At the order retained, the two sets of terms are additive provided that two conditions hold. The first is that the strain dependence of the nnn hoppings is neglected, which changes the tilt only by a relative amount of order $\tilde\gamma$ times the ratio of the nnn and nn hoppings, and this is small within the regime of validity stated below. The second is that the off-diagonal part is expanded around the Dirac point with the coefficients appropriate to the nn hopping. This second condition is not rigorous for a Dirac point displaced from $K$, because the numerical coefficients of the terms of order $a_0\,k$ in $d^\xi_x$ and $d^\xi_y$ then differ from those obtained by expanding around $K$, by amounts of the order of the displacement. We have not derived the corrected coefficients for a specific lattice. the Hamiltonian of Sec.~II~A is therefore adopted as a model for the coupling of tilted cones to strain, and we do not claim that it follows from a first-principles calculation for a particular material. The terms of order $a_0\,k_x$ matter for the results, since they are responsible for the intrinsic dispersion of the PLLs discussed in Sec.~II~B. The dispersion caused by the tilt does not depend on them. Statements on the lattice-induced part of the dispersion should be understood as statements of the model, and their material-specific values require the corresponding coefficients of the lattice in question.

The strain configuration is the following. Only the component $\varsigma_{yy}$ is nonzero, and it varies linearly along $y$. This describes a uniaxial strain along the direction of the bond $\bs M_1$ with a constant gradient along the same direction, which modulates the three nn hoppings in the ratios $\Delta u_1=4\,\Delta u_2=4\,\Delta u_3$. Other orientations, and shear components, generate different components of the pseudogauge field and are not covered by the Hamiltonian of Sec.~II~A. We restrict the tilt to the $k_x$-axis for a technical reason. For $\eta_y\neq0$, the term $v_y\,\eta_y\,k_y$ turns into $-i\,v_y\,\eta_y\,\partial_y$, and the equation for $\phi_B$ becomes a pair of coupled equations for its real and imaginary parts, which we could not solve analytically. Results for an arbitrary direction of the tilt are therefore not available in this work.

The expansion leading to the Hamiltonian of Sec.~II~A is controlled when $|\tilde\gamma|\ll1$ and $a_0\,|k_x|\ll1$. Since $\tilde\gamma=\gamma\,a_0\,\varsigma_{yy}$ and $\varsigma_{yy}=2\,b\,y/\gamma$, the first condition reads
\begin{align}
|\tilde\gamma|=\gamma\,a_0\,|\varsigma_{yy}|=2\,a_0\,e\,|\mathcal B|\,|y|\ll1\,,
\end{align}
and it restricts the strain to $|\varsigma_{yy}|\ll1/(\gamma\,a_0)$ in the region occupied by the states. For graphene, $\gamma\,a_0$ is of the order of 2 to 3, and strains of a few per cent keep $|\tilde\gamma|$ below $0.1$. For a strained region of width $W$, the PLL states must fit into the region, which requires $\ell_{\mathcal B}=1/\sqrt{e\,|\mathcal B|}\ll W$. The condition above, evaluated at $|y|\sim W$, requires $e\,|\mathcal B|\ll1/(2\,a_0\,W)$. Both conditions are fulfilled together if
\begin{align}
\frac{1}{W^2}\ll e\,|\mathcal B|\ll\frac{1}{2\,a_0\,W}\,.
\end{align}
For $W=100$ nm and $a_0=0.142$ nm, this window corresponds to $0.07\ {\rm T}\ll|\mathcal B|\ll23\ {\rm T}$ in conventional units. The condition $a_0\,|k_x|\ll1$ limits the part of the spectrum that can be trusted. The levels with large $n$, and the portions of each level with large $|k_x|$, lie beyond it. In particular, the finite upper bound on $n$ in the PLL spectrum indicates the loss of validity of the continuum model. For levels close to this bound, the spatial extent of the states reaches values for which $|\tilde\gamma|$ is of order unity. The results of this work should therefore be read as results for the low-lying PLLs, that is for $n$ well below the upper bound and for $a_0\,|k_x|\ll1$.''

%%%%%%%%%%%%%%%%%%%%%%%%%%%%%%%%%%%%%
\item \textcolor{blue}{Physical interpretation and validity of the pseudo-Landau levels. The authors note that ``a direct interpretation of the Hamiltonian as an analogy of actual LLs is a bit problematic because we have a mix of both $\Pi_x$ and $k_x$ terms''. If the PLLs are not dispersionless and do not arise from a genuine gauge field, the analogy is misleading. The authors should clarify what is the physical meaning of the PLLs and whether they are true bound states or a mathematical artefact of the hybrid spatial-momentum representation, why they remain dispersive, whether they represent localised cyclotron-like orbits or extended states, and what is the role of the tilting term $\eta_x\,k_x$, in particular whether it affects the localisation properties or the degeneracy. A clear physical picture and the conditions for experimental observation are needed.}

We agree that the physical picture was insufficiently developed. The PLLs are normalisable bound states of the strained Hamiltonian, and not an artefact of the hybrid representation. Translation invariance along $x$ is intact, and the $y$-dependence is solved exactly in the Appendix, where the eigenfunctions are generalised Laguerre polynomials times decaying factors, localised in $y$ over a length of the order of the pseudomagnetic length around a $k_x$-dependent centre. They are the analogue of cyclotron orbits whose guiding centre drifts along $x$. The dispersion has two sources. In a uniform real field, the translation of the guiding centre along $y$ is a symmetry, and the levels are flat. The strain does not have this symmetry, because the local hopping amplitudes and velocities vary with $y$, and this gives a dispersion of order $a_0\,k_x$. The tilt adds the term $v_x\,\eta_x\,k_x$. This term multiplies the identity matrix and therefore does not change the eigenfunctions, their localisation or the number of states per level, but it gives all levels the additional group velocity $v_x\,\eta_x$. The conditions of observability are discussed in the new Sec.~IV. The new text, placed at the end of Sec.~II~B, reads as follows.\\``The physical content of the PLL spectrum can be summarised as follows. For $\xi=+$ and small $\mathcal B$, the levels are
\begin{align}
\frac{E^{+}_{n,\rm si}}{v_x}\simeq\eta_x\,k_x+{\rm si}\,\sqrt{\left(1+\frac{3\,a_0\,k_x}{2}\right)\left(2\,|\zeta|\,n+e\,\mathcal B+|\zeta|\right)}\,.
\end{align}
At $k_x=0$ and for $n\gg1$, this reduces to the anisotropic Dirac ladder $\sqrt{2\,n\,v_x\,v_y\,e\,|\mathcal B|}$. Four ingredients are encoded in the full expression. First, the term $\eta_x\,k_x$ is a rigid tilt of all levels. It multiplies the identity matrix, and therefore changes neither the eigenfunctions nor the number of states per level, but it gives every level the group velocity $v_x\,\eta_x$ in addition to that of the untilted dispersion. Second, the factor $1+3\,a_0\,k_x/2$ stems from the strain-induced renormalisation of the velocities and from the terms of order $a_0\,k_x$ of the lattice, and it produces the dispersion of the untilted PLLs. Third, the valley index enters through $s_\xi$ and $r_\xi$, which is the origin of the valley asymmetry visible in Fig.~2. Fourth, the index $n$ has a finite upper bound, whose meaning is discussed in Sec.~II~A.

The PLLs are normalisable bound states, and not artefacts of the hybrid representation. Since the strain depends only on $y$, translation invariance along $x$ is intact and $k_x$ is a good quantum number, while the $y$-dependence is solved exactly in the Appendix. The eigenfunctions are generalised Laguerre polynomials multiplied by exponentially decaying factors. They vanish at the singular point $y=y^\xi_{\rm sg}$, where the coefficient of $k_y$ in $d^\xi_y$ vanishes, and they are localised in $y$ over a length of the order of $\ell_{\mathcal B}$ around a centre that depends on $k_x$. In the semiclassical picture, they correspond to cyclotron-like orbits whose guiding centre drifts along $x$ with the group velocity $v^{\xi,{\rm si}}_x(n;k_x)$. In a uniform real field, the energy of a level does not depend on $k_x$, because the translation of the guiding centre along $y$ is a symmetry. The strain does not have this symmetry. The local hopping amplitudes, and with them the local velocities, vary with $y$, and the energy of a state depends on where along $y$ it is located, that is on $k_x$. The PLLs are therefore dispersive even without tilt, and this part of the dispersion is a lattice effect of order $a_0\,k_x$. This is why the bulk longitudinal conductivity can be nonzero. At $T=0$, the electrical conductivity of a valley is
\begin{align}
\sigma^\xi_{xx}=\frac{e^2\,\tau}{L_y}\sum_{n,\,{\rm si}}\sum_{k_{\rm F}}\frac{|v^{\xi,{\rm si}}_x(n;k_{\rm F})|}{2\,\pi}\,,
\end{align}
where the sum over $k_{\rm F}$ runs over the Fermi points of each band. Every dispersive level that crosses the chemical potential contributes, while flat levels, for which $v_x=0$, do not.''

%%%%%%%%%%%%%%%%%%%%%%%%%%%
\item \textcolor{blue}{Semiclassical Boltzmann framework: applicability and limitations. A semiclassical Boltzmann treatment may not capture the quantum effects of Landau levels, such as phase coherence, level broadening and quantum oscillations, and its validity should be justified. The scattering mechanisms are not specified, although the transport coefficients depend critically on $\tau$, and it should be stated whether $\tau$ depends on energy, magnetic field or temperature. Is the relaxation-time approximation valid for the anisotropic and tilted system, and are more sophisticated approaches, such as variational methods, needed?}

The expression that we evaluate is the intraband (Fermi-surface) contribution computed with the exact PLL spectrum and velocities. The quantisation of the spectrum is therefore fully retained. What is semiclassical is the description of the scattering by a relaxation time and the neglect of the interband contributions, which are smaller by a factor of the order of $(\tau\,\Delta E)^{-2}$. The approach is therefore justified in the regime of resolved levels, $1/\tau\ll\Delta E$, and it does not apply when the disorder broadening merges the levels. The relaxation time is a phenomenological constant, independent of energy, momentum, temperature and $\mathcal B$, and the plotted quantities are normalised by it and their shapes therefore do not depend on its value. We have added a discussion of an energy-dependent $\tau$, of the limitations of the relaxation-time approximation for anisotropic scattering, and of the variational alternative, whose implementation we leave for future work since it requires the matrix elements of the impurity potential between PLL states. The new text, placed after the definitions of the coefficients in Sec.~III, reads as follows.\\``We comment on the status of this approach. the definitions of $\sigma_{xx}$, $\alpha_{xx}$ and $\ell_{xx}$ contains the intraband (Fermi-surface) contribution to the response, evaluated with the exact PLL spectrum and velocities. The effect of the quantisation on the spectrum is therefore fully retained. The semiclassical elements are the description of the scattering by a relaxation time and the neglect of the interband contributions, which involve the matrix elements of $v_x$ between different PLLs. With $\hbar=1$, these contributions are smaller than the intraband one by a factor of the order of $(\tau\,\Delta E)^{-2}$, where $\Delta E$ is the spacing between neighbouring PLLs. The approach is therefore justified when $1/\tau\ll\Delta E$ and $1/\tau\ll\mu$, which is the regime in which the PLLs are resolved. For $1/\tau\gtrsim\Delta E$, the disorder broadening merges the levels, the PLL features disappear, and the interband terms become important. Our results do not apply there. Phase-coherent effects, such as weak localisation, are not captured either.

The relaxation time is a phenomenological parameter, which we take to be independent of energy, momentum, temperature and $\mathcal B$. The plotted quantities $\tilde\sigma_{xx}$, $\tilde\alpha_{xx}$ and $\tilde\ell_{xx}$ are normalised by $\tau$, and their shapes do not depend on its value. An energy-dependent $\tau(\varepsilon)$ enters the integrals as a weight. At $T=0$, the conductivity then becomes $\sigma^\xi_{xx}=(e^2/L_y)\sum_{n,\,{\rm si}}\sum_{k_{\rm F}}\tau(\mu)\,|v_x|/(2\,\pi)$. The positions of the features are fixed by the band structure, while their heights are reweighted. The Mott relation and the Wiedemann-Franz law remain valid for any smooth $\tau(\varepsilon)$ as long as the scattering is elastic. A dependence of $\tau$ on the band index $n$ and on $k_x$ through the overlap of the PLL wave functions would modify the weights of the individual bands in the same manner. The relaxation-time approximation replaces the collision integral by $-\delta f/\tau$. For anisotropic and tilted spectra, and for strongly anisotropic scattering such as small-angle scattering, the transport time differs from the quantum lifetime. The solution of the linearised Boltzmann equation with the full collision integral, or a variational treatment [ziman], would then change the magnitudes by factors of order unity. It requires the matrix elements of the impurity potential between PLL states, which are overlaps of the Laguerre eigenfunctions of the Appendix, and it lies beyond the scope of this work.

Disorder and temperature act in a similar manner on the PLL features. The thermal broadening of the response is of the order of a few $k_{\rm B}\,T$, and the disorder broadening of the levels is of the order of $1/\tau$. The individual PLL features are resolved when both are smaller than $\Delta E$. The temperature dependence is displayed in Fig.~3, where the peaks broaden with decreasing $\tilde\beta$. The interband terms and the level broadening are not included in the Boltzmann treatment, and we do not study them in this work.''

%%%%%%%%%%%%%%%%
\item \textcolor{blue}{Transport coefficients: numerical results and interpretation. (a) The parameters used in the numerical calculations (temperature, chemical potential range, scattering time) are not clearly specified. (b) The observed discrepancies with the Mott relation and the Wiedemann-Franz law are significant, but no explanation is given of their origin, whether the energy dependence of the density of states, the PLLs or anisotropic scattering. (c) The authors do not explain how the oscillations in Fig.~6 depend on $\mathcal B$, temperature and chemical potential. (d) No specific experimental predictions are provided, such as expected oscillation frequencies and temperature dependence.}

We address the four points in turn. For point (a), all the plotted quantities are dimensionless. Energies are in units of $\hbar\,v_x/a_0$, and the temperature follows from $\tilde\beta$. The quantities $\tilde\sigma_{xx}$, $\tilde\alpha_{xx}$ and $\tilde\ell_{xx}$ are normalised by $\tau$. The scattering time therefore does not need to be specified for reproducibility. The remaining parameters are now stated in the captions of Figs.~3 to 6. For point (b), the explanation is given in the response to comment 5 below. The two Reviewers differ in their assessment of the agreement, and we find that both are consistent with the result that the quality of the agreement depends on the ratio of the temperature to the level spacing. For point (c), we have added the following paragraph, which gives the conversion to physical units and the positions and the scaling of the peaks. For point (d), the response to comment 9 below applies.\\``All the plotted quantities are dimensionless. Energies are measured in units of $\hbar\,v_x/a_0$, which is about $4.6$ eV for $v_x\simeq10^6$ m/s and $a_0=0.142$ nm, and the temperature follows from $k_{\rm B}\,T=(\hbar\,v_x/a_0)/\tilde\beta$. The quantities $\tilde\sigma_{xx}$, $\tilde\alpha_{xx}$ and $\tilde\ell_{xx}$ are normalised by $\tau$, as defined above, and therefore do not depend on its value. The remaining parameters of each figure are now stated in its caption. The positions of the peaks follow approximately from the band edges of the PLLs. In the absence of tilt, and for small $a_0\,|k_x|$, the peaks recur whenever $\mu^2/(2\,v_x\,v_y\,e\,|\mathcal B|)$ crosses an integer. Their spacing in $|\mathcal B|^{-1}$ is then $2\,v_x\,v_y\,e/\mu^2$, up to corrections from the tilt and from the terms of order $a_0\,k_x$. The peaks are resolved as long as the thermal broadening is smaller than the spacing of the levels, which grows as $\sqrt{|\mathcal B|}$.''


%%%%%%%%%%%%%%%%%%%%%5
\item \textcolor{blue}{Validity of the Mott relation and the Wiedemann-Franz law. The authors do not provide a rigorous justification for why these relations might or might not hold. The Mott relation is valid when the conductivity varies slowly with energy near the Fermi level, and the Wiedemann-Franz law assumes elastic scattering and a constant density of states, so that both assumptions are likely violated in the presence of PLLs. The authors should derive the conditions under which they hold, quantify the deviations as a function of temperature, magnetic field and chemical potential, and discuss the physical implications of these deviations for experimental measurements.}

We have derived the conditions of validity from the Sommerfeld expansion, and we have quantified the deviations. The derivation shows that, for the constant relaxation time used, the relations are controlled by the smoothness of the function $g(\varepsilon)$ on the scale of the temperature, and that they fail within about $T$ of the extrema of the PLL bands and of the end points of the branches, where $g$ has square-root singularities. We have computed the relative deviations for a range of temperatures and of the pseudomagnetic field, and for several values of the tilt, with the result that the Mott relation deviates by 15 to 20 per cent at temperatures below about a tenth of the level spacing and by less than 10 per cent beyond about 0.15 times the spacing, while the Wiedemann-Franz law holds at the level of a few per cent. The deviation at a given temperature grows with the pseudomagnetic field, and decreases with increasing tilt. The numbers were computed with the expressions of the manuscript for the stated parameters, and they are collected in a new table. The dependence on the chemical potential is the one described above, namely that the deviations are concentrated near the band edges, and the table refers to a window containing several of them. The new text, placed after the comparison with the Mott relation in Sec.~III, reads as follows. We also corrected the prefactor of the Mott relation, which contained $e^2$ instead of $e$ since $\alpha_{xx}$ carries a single factor of the charge. The figures use dimensionless variables and are not affected.\\``The relations in the Mott relation and the WF law follow from the Sommerfeld expansion of the three integrals in the definitions of $\sigma_{xx}$, $\alpha_{xx}$ and $\ell_{xx}$, which share the weight $-\partial_\varepsilon f_0\,(\partial_{k_x}\varepsilon)^2$. For an energy-independent $\tau$, they require no further input about the system. They hold when the function
\begin{align}
g(\varepsilon)=\frac{e^2\,\tau}{L_y}\sum_{n,\,{\rm si}}\int\frac{dk_x}{2\,\pi}\,\delta\left(\varepsilon-\varepsilon_n^{\rm si}\right)\left(\partial_{k_x}\varepsilon^{\rm si}_n\right)^2=\frac{e^2\,\tau}{L_y}\sum_{n,\,{\rm si}}\sum_{k_{\rm F}}\frac{|v_x|}{2\,\pi}
\end{align}
varies slowly on the scale of $T$ around $\mu$, with corrections of relative order $(\pi\,T/\delta)^2$, where $\delta$ is the energy scale on which $g$ varies. Here, $g$ is not smooth. Every extremum of a PLL band, and every end point of a branch, changes the number of Fermi points. Near a band minimum, $|v_x|$ grows as the square root of the distance in energy, and $\partial_\varepsilon g$ diverges there. In windows of width of the order of $T$ around these energies, the relations are not controlled. The Wiedemann-Franz law involves $g$ itself and not its derivative, and it is the more robust of the two. Both relations require elastic scattering, and they are exact in the Sommerfeld sense for the constant relaxation time used here. They are consistency checks of the smoothness of $g$, and they do not test the PLL physics itself.

We have quantified the deviations for $\xi=+$, $v_y/v_x=0.8$ and $\eta_x=0.3$ in the window $0.20\le\tilde\mu\le0.40$, restricted to $a_0\,|k_x|\le0.3$. With $\delta_{\rm M}$ and $\delta_{\rm WF}$ denoting the relative deviation, in the sense of the $L^2$ norm over the window, between the direct integral and the right-hand side of the Mott relation and of the Wiedemann-Franz law, respectively, the results are collected in Table~I. At fixed $\mathcal B$, the deviation of the Mott relation is of the order of 15 to 20 per cent for temperatures below about a tenth of the level spacing, and it falls below 10 per cent beyond about 0.15 times the spacing. At fixed temperature, it increases with $|\mathcal B|$, because the level spacing grows as $\sqrt{|\mathcal B|}$ and the relation is controlled by the ratio of $T$ to this spacing. Larger pseudomagnetic fields therefore break the Mott relation more strongly at a given temperature, within the range of validity of the continuum model. The Wiedemann-Franz law holds at the level of a few per cent throughout. A larger tilt reduces both deviations, for instance $\delta_{\rm M}=0.14$, $0.12$, $0.08$ and $0.06$ for $\eta_x=0$, $0.15$, $0.3$ and $0.45$ at $\tilde{\mathcal B}=0.01$ and $T=0.004$ (in units of $\hbar\,v_x/a_0$).
 For measurements, this implies that the thermopower at low temperatures cannot be obtained from the gate dependence of the conductivity through the Mott formula and has to be measured directly. The electronic part of the thermal conductivity can be inferred from the conductivity through the Wiedemann-Franz law to the accuracy stated here.''


%%%%%%%%%%%%%%%%%%%%%%%%%%%%%%%
\item \textcolor{blue}{Anisotropy and tilting: combined effects. The authors claim that their system acquires anisotropy from three ingredients, but the interplay between them is not systematically explored. How does the tilting angle affect the transport coefficients, given that $\eta_y=0$ is set and tilt is considered only along $k_x$? How does the ratio $v_y/v_x$, taken as $0.8$ in Fig.~6, affect the PLL spectrum and transport? Can the effects of strain and tilting be disentangled experimentally?}

The restriction to $\eta_y=0$ is technical. For $\eta_y\neq0$, the equation for $\phi_B$ becomes a pair of coupled equations for the real and imaginary parts, which we could not solve analytically, and this is now stated in Sec.~II~A (see the response to comment 1). The dependence on $\eta_x$ is explored in Fig.~4. The ratio $v_y/v_x$ enters the spectrum only through $\zeta=e\,\mathcal B\,v_y/v_x$, and we have added the corresponding scaling. We have also added a discussion of how the contributions of the strain and of the tilt can be separated. The new text, placed in Sec.~III after the discussion of the dependence on $\mathcal B$, reads as follows.\\``The dependence on the ratio $v_y/v_x$ enters the spectrum only through $\zeta=e\,\mathcal B\,v_y/v_x$. At leading order, a change of the ratio rescales the level spacing as $\sqrt{v_y/v_x}$ and the upper bound on $n$ as $v_x/v_y$, and it does not change the structure of the response. We have used $v_y/v_x=0.8$ as a representative value of moderate anisotropy. The effects of the strain and of the tilt can be separated in principle. The features that survive at $\eta_x=0$, which include the asymmetry of the response with respect to the sign of $\mathcal B$, are due to the strain, and their sign is reversed by reversing the strain gradient. The contributions $\propto\eta_x$ and $\propto\eta_x^2$ in Fig.~4 vanish at $\eta_x=0$, and the linear one is odd under the reversal of the tilt. They are therefore specific to the tilt, and they can be isolated by comparing measurements at two values of the tilt for the same strain profile.''

%%%%%%%%%%%%%%%%%
\item \textcolor{blue}{Comparison with existing literature. The comparison with the literature is superficial. The authors should clearly state what is novel compared with previous studies of PLLs in strained graphene, compare their results with those for untilted or isotropic Dirac cones, and discuss whether their findings are consistent with experimental observations of strain-induced pseudomagnetic fields in graphene.}

We have added a statement of the new ingredients to the ``Introduction'' section. They are the tilt-induced contribution to the group velocity, the thermoelectric and thermal coefficients, and the test of the Mott relation and the Wiedemann-Franz law. We have stated that the results reduce to those of the untilted case for $\eta_x=0$, which is the reference used in Fig.~4, and that the isotropic case corresponds to $v_y/v_x=1$ in the same formulas. We have also related the spectrum to the PLL sequences observed in strained graphene nanobubbles. The new text, placed at the end of the Introduction, reads as follows.\\``The new ingredients, compared with Refs.~[strained-dirac-franz, complex-LL-strained-graphene], are the tilt-induced contribution to the group velocity of the PLLs, the computation of the thermoelectric and thermal coefficients, and the test of the Mott relation and the Wiedemann-Franz law. For $\eta_x=0$, the spectrum reduces to the dispersive PLLs of the untilted anisotropic cone, and all the results below reduce to those of the untilted case. Scanning tunnelling spectroscopy of strained graphene nanobubbles [levy10] has resolved sequences of PLLs, which are consistent with the Dirac-like ladder that our spectrum reproduces at leading order [cf. Sec.~II~B]. The dispersion in $k_x$, which is the central ingredient for the longitudinal response, is not accessed by those experiments.''

%%%%%%%%%%%%%%%%%%%%%%%%
\item \textcolor{blue}{Presentation and clarity. Some equations are not clearly numbered or referenced in the text, for example Eq.~(6). Fig.~6 is not sufficiently explained, and the axis labels, units and parameter values should be clearly stated. The notation is sometimes inconsistent, for example $B$ versus $\mathcal B$ and $\zeta$ versus $\zeta_{yy}$.}

The equations are numbered automatically, and we have added the missing label for the group velocity and refer to the equations through their labels in the new text. The captions of Figs.~2 to 6 now state the definition of the plotted quantities, their normalisation and the parameters. Regarding the notation, $\mathcal B$ is the pseudomagnetic field and is the only magnetic field in the paper, while $\varsigma_{yy}$ is the strain component and $\zeta=e\,\mathcal B\,v_y/v_x$ is a different quantity. The similarity of the two symbols was a source of confusion, and we have added the following paragraph after the definition of the strain configuration, which fixes the notation. We have also corrected the prefactor of the Mott relation, as mentioned in the response to comment 5.\\``Notation: the symbol $\varsigma_{ij}$ denotes the strain tensor, of which only $\varsigma_{yy}$ is nonzero. The pseudomagnetic field is denoted by $\mathcal B$ throughout, and a real magnetic field does not appear in this work. We also use $b=e\,\mathcal B$ and $\zeta=e\,\mathcal B\,v_y/v_x$. The latter must not be confused with $\varsigma_{yy}$. The dimensionless variables carrying a tilde are defined in Sec.~III.''

\item \textcolor{blue}{Experimental relevance and predictions. The claim of ``unambiguous experimental signatures'' is not substantiated. The authors should propose specific experimental setups, such as strained graphene on a substrate or bent nanoribbons, estimate measurable quantities, such as conductivity oscillations and thermoelectric voltage, and compare them with typical experimental resolutions, and discuss the challenges of realising the required strain configurations.}

We agree that the claim was not supported, and we have removed the word ``unambiguous'' from the abstract. We have added a new section, Sec.~IV, with the candidate platforms, the observables, estimates of the scales and the practical challenges. It reads as follows.\\``\section{Experimental considerations}
The calculation is most directly relevant for platforms in which tilted Dirac cones and controllable strain gradients coexist. Strained graphene is the platform in which pseudomagnetic fields have been demonstrated, with pseudo-fields of several hundred tesla in nanobubbles [levy10, graphene-strain-geim, vozmediano10, strain-2d]. A tilt can be induced in the deformed lattice, as discussed in Ref.~[tilted-goerbig]. Organic conductors of the $\alpha$-(BEDT-TTF)$_2$I$_3$ family under pressure host tilted Dirac cones with a large tilt [organic-tilted-dirac], and the tilt is tunable by the pressure and by its direction. Two-dimensional borophene with tilted anisotropic cones [zhou14] is a further candidate, in which the strain response has still to be explored.

The observables are the following. The first is the longitudinal conductivity as a function of the chemical potential, tuned by a gate, whose peak structure follows the band edges of the PLLs. The second is the dependence of the conductivity on the tilt, $\delta\sigma$ of Fig.~4, which has a term linear in $\eta_x$ for $\mu\neq0$ and no such term at $\mu=0$. The third is the asymmetry with respect to the sign of the pseudomagnetic field, which can be probed by reversing the strain gradient. The fourth is the thermopower $S=\alpha_{xx}/\sigma_{xx}$. By the Mott relation, $S\sim(\pi^2/3)\,(k_{\rm B}/e)\,(k_{\rm B}\,T/\delta)$, which is a few tens of $\mu$V/K for $k_{\rm B}\,T/\delta\sim0.1$ and is within the range of thermopower measurements in graphene.

We estimate the scales as follows. Each Fermi point contributes $e^2\,v\,\tau/(h\,L_y)$ to the conductivity of a valley, that is $e^2/h$ times the ratio of the mean free path to the width. For a mean free path comparable to the width, the contribution of each Fermi point is of the order of $e^2/h\simeq39\ \mu$S, which is well above the resolution of standard transport measurements. The level spacing at $k_x=0$ follows from the Dirac-like ladder, $E_n\simeq\sqrt{2\,n\,v_x\,v_y\,e\,|\mathcal B|}$. For $v_x\simeq v_y\simeq10^6$ m/s, this is about $36\ {\rm meV}\,\sqrt{n\,|\mathcal B|/{\rm T}}$. The PLLs are then resolved at temperatures of a few tens of kelvin for a pseudo-field of about $1$ T. In organic conductors, where the velocities are about an order of magnitude smaller, the spacings are reduced by the same factor. The pseudo-field is related to the strain gradient by $e\,\mathcal B=(\gamma/2)\,\partial_y\varsigma_{yy}$. For $\gamma\,a_0\simeq3$, a strain variation of $1\%$ over $100$ nm corresponds to $|\mathcal B|\simeq0.7$ T. PLL spectroscopy at liquid-helium temperatures [levy10] shows that the disorder broadening in graphene can remain below the spacing of the levels, which is the condition $1/\tau\ll\Delta E$ for the validity of our treatment.

There are practical challenges. The first is the realisation of a strain gradient that is constant over the width of the sample, since the pseudo-field varies if the gradient is not constant. The second is the tuning of the tilt independently of the strain, which in organic conductors is achieved by pressure and which also changes the velocities. The third concerns the valleys. A transport measurement yields the sum of the contributions of both valleys, whereas our figures refer to the valley $\xi=+$. The contribution of the valley $\xi=-$ follows from the PLL spectrum for $\xi=-$, and it has a different range of the PLLs owing to the strain-induced shift of the Dirac points. The two contributions therefore do not cancel. We do not propose a valley-resolved transport measurement, and the comparison with experiment should be made for the sum. Finally, the strain also shifts the chemical potential, and a gate is needed to separate this effect from that of the PLLs.''

\item \textcolor{blue}{Lack of discussion on disorder and temperature effects. The authors should include a phenomenological scattering rate and study its effect on the transport coefficients, discuss the temperature dependence of the PLL features and whether they survive at finite temperature, and compare with experiments in which disorder and temperature effects are known to be significant.}

The phenomenological scattering rate is the relaxation time $\tau$, which multiplies all the coefficients and does not change their shapes within the approach. The broadening of the levels by disorder is not contained in the Boltzmann treatment, and we have stated the criterion $1/\tau\ll\Delta E$ for the validity of the results together with the statement that thermal and disorder broadening act in a similar manner on the PLL features. The temperature dependence is shown in Fig.~3, and the survival of the features is quantified in the new text in terms of the ratio of a few $k_{\rm B}\,T$ to the level spacing. The comparison with experiments is made in Sec.~IV through the observation that PLLs are resolved in scanning tunnelling spectroscopy of strained graphene at liquid-helium temperatures. The relevant paragraphs are those quoted in the responses to comments 3 and 9.
\end{enumerate}


%%%%%%%%%%%%%%%%%%%%%%%%
\section{Response to Reviewer \#3}

\begin{enumerate}

\item \textcolor{blue}{The physical distinction between pseudo-Landau levels and true Landau levels needs a clearer physical discussion. The authors emphasise that the PLLs are dispersive with a finite longitudinal group velocity, unlike the flat conventional Landau levels. However, the physical origin of the valley-dependent pseudogauge field and its key differences from real magnetic fields should be more explicitly contrasted in the main text. In particular, pseudomagnetic fields do not break time-reversal symmetry, act oppositely on the two valleys, and do not produce a true quantum Hall effect. At several points the text draws analogies to Landau quantisation without warning that many standard quantum-Hall phenomena cannot be imported. The discussion should elaborate why the bulk longitudinal conductivity can be nonzero for dispersive PLLs, which is a central result of the paper.}

We thank the Referee for this comment, and we agree with the points listed. We have added a paragraph to the Introduction that contrasts the pseudomagnetic field with a real field and warns that the quantum-Hall phenomenology does not carry over. The explanation of why the bulk longitudinal conductivity is nonzero is given in the new text of Sec.~II~B, which is quoted in the response to comment 2 of Reviewer \#1, and it is based on the expression of $\sigma_{xx}$ as a sum over the Fermi points of the dispersive bands. The new paragraph of the Introduction reads as follows.\\``The pseudomagnetic field must not be confused with a real magnetic field, and several results of the standard quantum-Hall phenomenology do not carry over to it. The strain-induced gauge field acts with opposite signs on the two valleys. Time reversal maps one valley onto the other and, since it also reverses the momentum, it leaves the strain-induced field invariant. A strain gradient therefore preserves time-reversal symmetry, in contrast to a real magnetic field, which breaks it. As a consequence, the Hall conductivity summed over the two valleys must vanish, and neither quantised Hall plateaux nor a net chiral edge current are expected from the bulk response studied here. At most, a valley-Hall type response with opposite signs in the two valleys can occur. The longitudinal contributions of the two valleys, on the other hand, do not cancel. Conventional Landau levels are flat in momentum, because a uniform field is invariant under translations along $y$ combined with a gauge transformation, which fixes the guiding centre as $y_0=-k_x/(e\,B)$ without any dependence of the energy on it. Flat bands carry no current in the Boltzmann picture and the bulk longitudinal conductivity of a clean system then vanishes. In the strained system, this invariance is absent, because the strain also renormalises the hopping amplitudes, and hence the local Fermi velocities, as functions of $y$. The tilt adds a further source of dispersion. The origin of the dispersion and its consequence, a nonzero bulk longitudinal conductivity, are discussed in Sec.~II~B.''

%%%%%%%%%%%%%%%%%
\item \textcolor{blue}{Approximation validity and model assumptions need a more critical assessment. The work adopts a semiclassical Boltzmann formalism with a momentum-independent relaxation time $\tau$. The authors should discuss when this approximation is justified, how an energy- or momentum-dependent scattering would modify the transport coefficients, and how disorder could smear the PLL-related peaks. The authors assume that the strain formalism derived for the $K/K^\prime$ points of graphene carries over to tilted Dirac cones with Dirac points away from the high-symmetry points. More careful validation or caveats are required. How large can the strain be before the continuum low-energy Hamiltonian breaks down, and what are the bounds on the strain magnitude?}

The discussion of the validity of the relaxation-time approach, of the effect of an energy-dependent or band-dependent scattering time, and of the smearing of the peaks by disorder has been added to Sec.~III, and the text is quoted in the response to comment 3 of Reviewer \#1. The caveats on the extension of the strain formalism to tilted cones and the bounds on the strain are in the new text of Sec.~II~A, which is quoted in the response to comment 1 of Reviewer \#1. In summary, the continuum Hamiltonian is controlled when $|\tilde\gamma|=\gamma\,a_0\,|\varsigma_{yy}|\ll1$ in the region occupied by the states, which is a strain of a few per cent for graphene, and when $a_0\,|k_x|\ll1$. For a strained region of width $W$, the pseudomagnetic field must satisfy $1/W^2\ll e\,|\mathcal B|\ll1/(2\,a_0\,W)$. We have also stated that the finite upper bound on the PLL index indicates the loss of validity of the model.\item \textcolor{blue}{The Mott relation and Wiedemann-Franz law results require a deeper physical interpretation. The numerical comparison shows reasonably good agreement between the direct integral computation and the Mott and Wiedemann-Franz relations, even for this PLL-hosting system. However, the manuscript only states that ``the relations produce a pretty good fit'' without explaining why. These relations are derived assuming a smooth energy dependence of the conductivity near the Fermi energy, while the PLL spectrum consists of discrete levels. The authors should explain under what thermal smearing conditions the relations remain approximately valid, point out the visible small deviations, and comment on whether larger strain or larger pseudomagnetic fields would break the relations more strongly.}

The new text quoted in the response to comment 5 of Reviewer \#1 gives the conditions under which the relations hold and identifies where the deviations are located. In particular, the deviations are concentrated within about $T$ of the extrema of the PLL bands and of the end points of the branches. We have also examined the question about larger fields. Within the range of validity of the continuum model, a larger pseudomagnetic field breaks the Mott relation more strongly at a given temperature, because the level spacing grows as $\sqrt{|\mathcal B|}$ and the relation is controlled by the ratio of the temperature to this spacing. This is quantified in the new table. The Wiedemann-Franz law is more robust.\item \textcolor{blue}{Experimental feasibility and observable signatures should be expanded. The abstract mentions ``unambiguous experimental signatures'', but the discussion of the experimental realisation is very brief. Which material platforms are most promising, for example organic conductors under pressure or strained two-dimensional van der Waals materials? Which concrete observables can experimentalists target, such as the peak structure of the longitudinal conductivity as a function of the chemical potential, the tilt-parameter dependence of the conductivity ($\delta\sigma$ in Fig.~4) and the asymmetry with respect to the sign of the pseudomagnetic field? The authors should discuss practical challenges, namely how to tune the tilt $\eta_x$ controllably, how to generate the required uniaxial strain profile, and how to separate the contributions of the two valleys in a measurement.}

The new Sec.~IV, quoted in full in the response to comment 9 of Reviewer \#1, discusses the platforms (strained graphene, organic conductors under pressure and borophene), the three observables named by the Referee together with the thermopower, the estimates of the scales, and the challenges of tuning the tilt, of producing the strain gradient and of the valley contributions. Regarding the valleys, a transport measurement yields the sum of the two valley contributions, and we do not propose a valley-resolved measurement. We have stated that the comparison with experiment should be made for the sum. The word ``unambiguous'' has been removed from the abstract.\item \textcolor{blue}{Presentation of the PLL eigenvalue formula and mathematical complexity. The energy dispersion formula for the PLLs is cumbersome. The small-$\mathcal B$ approximate expression is given, but physical intuition for the full expression is limited. The main text should highlight the key physical effects encoded in the formula, namely the tilt-induced shift $\eta_x\,k_x$, the dispersive $k_x$ dependence, the valley asymmetry and the finite upper bound on the level index $n$. The key physical takeaways should be separated from the dense mathematical expressions, and some intermediate steps of the appendix could gain short physical comments.}

We thank the Referee for this suggestion. The first paragraph of the new text in Sec.~II~B, quoted in the response to comment 2 of Reviewer \#1, gives the small-$\mathcal B$ expression for the valley $\xi=+$ and explains in words the four ingredients named by the Referee. In the Appendix, we have added two short physical comments. The first, after the asymptotic analysis at $y\to0$, states that the admissible solution is the one that vanishes at the point where the strain-reduced velocity along $y$ vanishes. The second, after the identification of the Laguerre polynomials, states that the truncation of the confluent hypergeometric series is the analogue of the Landau quantisation condition, that $n$ counts the nodes of the wave function along $y$, and that the dependence of $\sqrt\Delta$ on $k_x$ is what makes the levels dispersive.
\end{enumerate}

%%%%%%%%%%%%%%%%%%%%%%%%%%%%%%%%%%%%%%%%%%%%%%
\section{List of changes}

All the additions are marked in the manuscript with the colour red. The word ``unambiguous'' has been removed from the abstract. In the Introduction, a paragraph contrasting the pseudomagnetic field with a real field and a statement of the new ingredients of the work have been added. In Sec.~II~A, the sentence presenting the extension of the strain coupling to tilted cones as justified has been replaced by a motivated assumption, the assumptions, the strain configuration and the conditions of validity have been stated, the reason for setting $\eta_y=0$ has been given, and the notation has been fixed. In Sec.~II~B, a summary of the physical content of the PLL spectrum and a discussion of the nature of the PLLs and of the origin of the nonzero longitudinal conductivity have been added. In Sec.~III, the status of the Boltzmann approach, the role of the scattering time, and the effects of disorder and temperature have been discussed, the parameters and units of the figures have been stated in the captions, the scaling of the peaks has been given, the dependence on $v_y/v_x$ and the separation of the strain and tilt effects have been discussed, and the conditions of validity of the Mott relation and the Wiedemann-Franz law have been derived and the deviations quantified in a new table. The prefactor $e^2$ of the Mott relation has been corrected to $e$, and this does not affect any figure. A new Sec.~IV on the experimental considerations has been added. Two short physical comments have been added to the Appendix. 

\end{document}
